\begin{lemma}
\label{lemma-iso-sheaf}
In Situation \ref{situation-iso}.
\begin{enumerate}
\item Each of the functors $F_{iso}$, $F_{inj}$, $F_{surj}$, $F_{zero}$
satisfies the sheaf property for the fpqc topology.
\item If $f$ is quasi-compact and $\mathcal{G}$ is of finite type,
then $F_{surj}$ is limit preserving.
\item If $f$ is quasi-compact and $\mathcal{F}$ of finite type, then
$F_{zero}$ is limit preserving.
\item If $f$ is quasi-compact, $\mathcal{F}$ is of finite type, and
$\mathcal{G}$ is of finite presentation, then $F_{iso}$ is limit preserving.
\end{enumerate}
\end{lemma}

\begin{proof}
Let $\{T_i \to T\}_{i \in I}$ be an fpqc covering of schemes over $S$.
Set $X_i = X_{T_i} = X \times_S T_i$ and $u_i = u_{T_i}$.
Note that $\{X_i \to X_T\}_{i \in I}$ is an fpqc covering of $X_T$, see
Topologies, Lemma \ref{topologies-lemma-fpqc}.
In particular, for every $x \in X_T$ there exists an $i \in I$
and an $x_i \in X_i$ mapping to $x$. Since
$\mathcal{O}_{X_T, x} \to \mathcal{O}_{X_i, x_i}$ is flat, hence
faithfully flat (see
Algebra, Lemma \ref{algebra-lemma-local-flat-ff})
we conclude that $(u_i)_{x_i}$ is injective, surjective, bijective, or zero
if and only if $(u_T)_x$ is injective, surjective, bijective, or zero.
Whence part (1) of the lemma.

\medskip\noindent
Proof of (2). Assume $f$ quasi-compact and $\mathcal{G}$ of finite type.
Let $T = \lim_{i \in I} T_i$ be a directed limit of affine $S$-schemes
and assume that $u_T$ is surjective.
Set $X_i = X_{T_i} = X \times_S T_i$ and
$u_i = u_{T_i} : \mathcal{F}_i = \mathcal{F}_{T_i}
\to \mathcal{G}_i = \mathcal{G}_{T_i}$.
To prove part (2) we have to show that $u_i$ is surjective for some $i$.
Pick $i_0 \in I$ and replace $I$ by $\{i \mid i \geq i_0\}$.
Since $f$ is quasi-compact the scheme $X_{i_0}$ is quasi-compact.
Hence we may choose affine opens $W_1, \ldots, W_m \subset X$
and an affine open covering
$X_{i_0} = U_{1, i_0} \cup \ldots \cup U_{m, i_0}$ such that
$U_{j, i_0}$ maps into $W_j$ under the projection morphism $X_{i_0} \to X$.
For any $i \in I$ let $U_{j, i}$ be the inverse image of $U_{j, i_0}$.
Setting $U_j = \lim_i U_{j, i}$ we see that $X_T = U_1 \cup \ldots \cup U_m$
is an affine open covering of $X_T$. Now it suffices to show, for a given
$j \in \{1, \ldots, m\}$ that $u_i|_{U_{j, i}}$ is surjective for some
$i = i(j) \in I$. Using
Properties, Lemma \ref{properties-lemma-finite-type-module}
this translates into the following algebra problem:
Let $A$ be a ring and let $u : M \to N$ be an $A$-module map.
Suppose that $R = \colim_{i \in I} R_i$ is a directed colimit
of $A$-algebras. If $N$ is a finite $A$-module and if
$u \otimes 1 : M \otimes_A R \to N \otimes_A R$ is surjective, then
for some $i$ the map
$u \otimes 1 : M \otimes_A R_i \to N \otimes_A R_i$ is surjective.
This is
Algebra, Lemma \ref{algebra-lemma-module-map-property-in-colimit} part (2).

\medskip\noindent
Proof of (3). Exactly the same arguments as given in the proof of (2)
reduces this to the following algebra problem:
Let $A$ be a ring and let $u : M \to N$ be an $A$-module map.
Suppose that $R = \colim_{i \in I} R_i$ is a directed colimit
of $A$-algebras. If $M$ is a finite $A$-module and if
$u \otimes 1 : M \otimes_A R \to N \otimes_A R$ is zero, then
for some $i$ the map
$u \otimes 1 : M \otimes_A R_i \to N \otimes_A R_i$ is zero.
This is
Algebra, Lemma \ref{algebra-lemma-module-map-property-in-colimit} part (1).

\medskip\noindent
Proof of (4).
Assume $f$ quasi-compact and $\mathcal{F}, \mathcal{G}$ of finite presentation.
Arguing in exactly the same manner as in the previous paragraph
(using in addition also
Properties, Lemma \ref{properties-lemma-finite-presentation-module})
part (3) translates into the following algebra statement:
Let $A$ be a ring and let $u : M \to N$ be an $A$-module map.
Suppose that $R = \colim_{i \in I} R_i$ is a directed colimit
of $A$-algebras. Assume $M$ is a finite $A$-module, $N$ is a finitely
presented $A$-module, and
$u \otimes 1 : M \otimes_A R \to N \otimes_A R$ is an isomorphism.
Then for some $i$ the map
$u \otimes 1 : M \otimes_A R_i \to N \otimes_A R_i$ is an isomorphism.
This is
Algebra, Lemma \ref{algebra-lemma-module-map-property-in-colimit} part (3).
\end{proof}
